%% notes-tex/031-unified-representation/sections/3-established.tex

\section{What the data exploration established\secstatus{tent}}
\label{sec:established}

The earlier document, \texttt{notes-tex/031-inhibitor-tolerance-data}, compared Vanacloig
with the two Hillenmeyer arms axis by axis. Three of its results bound everything planned
here, and the linear and nearest-neighbor baselines run since then say how much of the bound
a simple model reaches.

\notesec{Replication sets the ceiling, and the served error is honest for one dataset only}
The reliability index of a condition is one minus the mean squared served standard error over
the variance of the response, and its square root is the ceiling on the correlation any model
can reach with the noise-free response. Table~\ref{tab:reliability} places that index beside
the raw replicate agreement of the same conditions and beside what the Spearman-Brown lift of
the single-replicate reliability predicts for the mean of $k$ replicates. For Vanacloig the
index sits 0.04 above the prediction and its median ceiling is 0.84, with three of 41
conditions carrying no signal at all, 5-hydroxymethylfurfural among them. For both Hillenmeyer
arms the index sits 0.28 and 0.36 above the prediction, so their served error understates the
noise and the raw replicate agreement, 0.39 and 0.26, is the number to plan against. Hoepfner
and Wildenhain serve no uncertainty on most records, so no ceiling exists for them yet.

%% SOURCE: experiments/031-env-chemgen-inhibitor-tolerance/scripts/replicate_noise_and_ceilings.py
%% and reliability_reconciliation.py -> results/noise_summary.csv and
%% results/reliability_reconciliation.csv, rendered by notes_tex_representation_tables.py.
\begin{table}[htbp]
\centering
\footnotesize
\caption{Reliability per dataset, medians over conditions. Replicate rho is the raw
between-replicate Spearman over genes; rel.\ single is the single-replicate reliability it
implies; $k$ the replicate count; predicted is the Spearman-Brown reliability of the mean of
$k$; index is the served-error reliability; ceiling is $\sqrt{\text{index}}$; no signal counts
conditions whose index is at or below zero.}
\label{tab:reliability}
%% GENERATED by experiments/031-env-chemgen-inhibitor-tolerance/scripts/notes_tex_representation_tables.py
%% SOURCE: experiments/031-env-chemgen-inhibitor-tolerance/results/noise_summary.csv and reliability_reconciliation.csv from replicate_noise_and_ceilings.py and reliability_reconciliation.py
%% Do not edit by hand.
\begin{tabular}{lrrrrrrrrr}
\toprule
\textbf{dataset} & \textbf{conditions} & \textbf{replicate rho} & \textbf{rel. single} & \textbf{k} & \textbf{predicted} & \textbf{index} & \textbf{index minus predicted} & \textbf{ceiling} & \textbf{no signal} \\
\midrule
Vanacloig 2022 & 41 & 0.33 & 0.39 & 3 & 0.66 & 0.70 & +0.04 & 0.84 & 3 \\
Hillenmeyer HOM & 140 & 0.39 & 0.42 & 4 & 0.59 & 0.92 & +0.28 & 0.96 & 0 \\
Hillenmeyer HET & 308 & 0.26 & 0.25 & 4 & 0.42 & 0.84 & +0.36 & 0.92 & 2 \\
\bottomrule
\end{tabular}

\end{table}

%% SOURCE: experiments/031-env-chemgen-inhibitor-tolerance/scripts/replicate_noise_and_ceilings.py ->
%% vanacloig_ceilings.svg, from results/condition_noise_vanacloig2022.csv, converted by `make plots`.
\tcfig{figures/vanacloig_ceilings.pdf}{Every Vanacloig compound, ordered by its ceiling on the
correlation with the noise-free response (bars, $\sqrt{\mathrm{rel}}$), with the raw three-batch
replicate Spearman (circles) and the between-batch hit Jaccard (squares). A compound with no
bar has an index at or below zero.}{fig:ceilings}

\notesec{Cross-dataset transfer is a gene prior, not compound-matched transfer}
Between Vanacloig and the Hillenmeyer homozygous arm the per-gene mean response correlates at
0.118 and the gene-gene similarity structure agrees eleven null standard deviations from zero,
while the median Spearman between matched condition profiles is 0.010. The attenuation limit
set by the two reliabilities is 0.55, so the failed transfer is not a noise artifact. Exactly
one compound, methyl methanesulfonate, has a per-gene profile that transfers; benomyl, the
acetate salt against the free acid, and ferulic acid at two doses do not. Chemical similarity
correlates with response similarity across the two datasets at Spearman 0.09 at best
(Section~\ref{sec:encoders}).

\notesec{Within Vanacloig, molecule features carry real signal and gene embeddings do not}
Table~\ref{tab:baselines} gives the best ridge and nearest-neighbor model per split and
target beside the nulls. On the compound cold-start split, leaving one of the 41 compounds out
in turn, the centered target subtracts each gene's mean over the training compounds so that
only compound-specific signal remains. There a ridge map from the FCFP4 fingerprint reaches
Spearman 0.311 against 0.020 for a random training neighbor, positive in 93 percent of folds,
which is 42 percent of the 0.835 ceiling. On the raw target the gene mean alone reaches
0.212 and the best ridge 0.341. On the gene cold-start split, five gene folds by 41 compounds,
protein-language-model embeddings reach 0.044, real but 8 percent of the ceiling. Ridge beats
every nearest-neighbor variant, and the encoder ranking does not favor the large pretrained
models: a plain feature-class fingerprint is best on the centered target.

%% SOURCE: experiments/031-env-chemgen-inhibitor-tolerance/scripts/baseline_ceilings.py ->
%% results/baseline_ceilings_summary.csv, rendered by notes_tex_representation_tables.py.
\begin{table}[htbp]
\centering
\footnotesize
\caption{Linear and nearest-neighbor baselines on Vanacloig against the reliability ceiling,
medians over folds. For each split and target the best feature set per model class is shown.
Compound cold holds out one compound per fold; gene cold holds out a fifth of the genes per
fold and scores each compound column. The centered target removes each gene's training mean.}
\label{tab:baselines}
%% GENERATED by experiments/031-env-chemgen-inhibitor-tolerance/scripts/notes_tex_representation_tables.py
%% SOURCE: experiments/031-env-chemgen-inhibitor-tolerance/results/baseline_ceilings_summary.csv from baseline_ceilings.py
%% Do not edit by hand.
\begin{tabular}{llllrrrrr}
\toprule
\textbf{split} & \textbf{target} & \textbf{model} & \textbf{features} & \textbf{folds} & \textbf{Spearman} & \textbf{Pearson} & \textbf{ceiling} & \textbf{of ceiling} \\
\midrule
compound cold & centered & ridge & fcfp4 count & 41 & 0.311 & 0.350 & 0.835 & 0.42 \\
compound cold & centered & knn5 & unimol v1 & 41 & 0.284 & 0.328 & 0.835 & 0.39 \\
compound cold & centered & knn3 & fcfp4 count & 41 & 0.283 & 0.327 & 0.835 & 0.39 \\
compound cold & centered & knn1 & molformer xl & 41 & 0.280 & 0.279 & 0.835 & 0.33 \\
compound cold & centered & mean train profile & none & 41 & -0.001 & -0.011 & 0.835 & -0.01 \\
compound cold & raw & ridge & mole static & 41 & 0.341 & 0.393 & 0.838 & 0.47 \\
compound cold & raw & knn5 & chemberta2 mlm & 41 & 0.312 & 0.401 & 0.838 & 0.48 \\
compound cold & raw & knn3 & roberta zinc 480m & 41 & 0.287 & 0.367 & 0.838 & 0.44 \\
compound cold & raw & knn1 & fcfp4 count & 41 & 0.281 & 0.338 & 0.838 & 0.40 \\
compound cold & raw & mean train profile & none & 41 & 0.212 & 0.255 & 0.838 & 0.30 \\
compound cold & raw & gene mean & none & 41 & 0.212 & 0.255 & 0.838 & 0.30 \\
gene cold & centered & ridge & prott5 & 205 & 0.044 & 0.068 & 0.841 & 0.08 \\
gene cold & raw & ridge & prott5 & 205 & 0.044 & 0.068 & 0.841 & 0.08 \\
\bottomrule
\end{tabular}

\end{table}

Two consequences carry into the plan. More than half the predictable variance on the
compound-cold target is unclaimed by a linear map, so there is headroom for a model that uses
the gene graph. And within-dataset chemistry works far better than cross-dataset chemistry,
0.31 against at most 0.09, which is what one medium, one dose basis and one readout buy; the
dataset token and the dose basis token in Section~\ref{sec:representation} are how the pooled
model is told which regime a record came from.
